本研究は、特許文書を用いたクロスドメイン技術候補抽出において、全文類似だけでは候補抽出理由を説明しにくいという問題に着目する。特許文書には、技術課題、解決手段、対象物、効果などが含まれるが、全文ベクトルとして一括で扱うと、どの観点で候補が近いのかが不明瞭になりやすい。そこで本研究では、特許文書を課題文脈である \texttt{problem\_context} と解決手段文脈である \texttt{solution\_context} に分離し、課題が近い異分野間において、対象分野では低出現または未使用である技術手法・解決手段を専門家確認候補として抽出する枠組みを検討した。実験1では、U-Netを対象とした近接解決策型の主分析を行い、医療画像解析分野で先行して現れたU-Netを、製造欠陥検出分野での後年出現を参照しながら後ろ向き評価した。その結果、U-Netは過去期間の具体的手法ランキングで2位となり、上位の専門家確認候補として確認された。実験2では、洗浄・除去系技術から半導体洗浄への異質解決策型の補助分析を行った。46件のユニークファミリーをLLM支援付き著者確認評価により確認した結果、◎11件、○5件、△12件、×18件となった。ノイズも含まれた一方で、一部クラスタでは専門家確認候補として解釈可能な候補が得られた。また、軽量全文テキストベースラインとの比較により、提案手法と全文テキストベースラインの候補集合に差分があることを補助的に確認した。ただし、本研究は技術導入可能性を証明するものではなく、専門家確認候補を絞り込むスクリーニング手法として位置づける。
